% Y. Seirawan - G. Kasparov, Skelleftea World Cup 1989: your comments. Format: see src/endgame_book/gamebook.py.
@ title | Y. Seirawan - G. Kasparov, Skelleftea World Cup 1989

@ 11.f3 | board
\mythoughts{I am not sure about the result after 11...h5: my g5 square is weakened. Without knowledge of
the opening, I will play 11...\K h7.}

@ 11...h5
\myreflection{``The daring pawn thrust is the modern method of active defense, which has successfully
passed multiple practical trials, both in Indian systems and in the Sicilian Dragon.''
\cite{Nikitin:2020-coaching-kasparov-2}}

@ 12.a4
\mythoughts{A typical plan for Black in the Benoni is ...a6, ...b5 and ...c4, and moving the b8 knight to
d3 via d7 and c5. White's typical ideas are attacking the d6-pawn with a knight on c4, or the e5 thrust.
He can do neither of these, so my position looks comfortable.}

@ 13.a5
\mythoughts{13.a5 is a cold shower: the previous plan is no longer possible. I must have a plan to develop
my queenside pieces. My bishop has nowhere to go, since ...\B d7 would take the square my knight needs.
I want to bring my knight to e5 via d7. As long as White cannot play f4, my knight is safe there.
I want to play 13...\N bd7.}

``A questionable decision. If White fails to take advantage of the outpost on b6, the pawn's raid
will not only turn out to have been a waste of time, but will also force White to constantly tend
to its protection.'' \cite{Nikitin:2020-coaching-kasparov-2}

@ 13...Nh7
\myreflection{My 13...\N bd7 is as good as the text move. The text move 13...\N h7 is more dynamic
and allows active play such as ...g5 on the kingside.}

@ 14.Bd1 | board
\mythoughts{What does White want? First, to free the e2 square for the development of his knight.
Second, \B c2, protecting the e4-pawn.

I want to play 14...f5. After 14...f5 15.exf5 \B h6 16.\K f2 \B xf5, Black looks good.

15.\B c2 fxe4 16.fxe4 \B g4: Black is not afraid of h3, because of ...\Q h4+.}

@ 14...Nd7
\myreflection{Asking what White wants is the only good part of my thinking here; everything else is
nonsense. 14...f5 does not hinder White's plan at all, and my calculation is flawed:

14...f5 15.exf5? \B h6 16.\K f2 \B xf5? leaves the bishop on h6 hanging. 16...\B xe3+, winning material,
is the right move.

14...f5 15.\B c2 fxe4 16.fxe4 \B g4? ``Black is not afraid of h3, because of ...\Q h4+.'' But 16...\B g4 does
nothing, and 17.h3 \Q h4+ is answered by 18.\B f2. Black is simply losing.

I must check again and again whether my logic is sound and whether it breaks the laws of chess. Here the
logic is not there. The move also breaks a law of chess: \term{do-not-weaken}.\index{Strategy!Do not weaken yourself}
I weaken my pawn structure prematurely, without a reason. Such a move should not be calculated at all,
let alone calculated wrongly.}

@ 16.b3
\mythoughts{I want to make White's development as difficult as possible. 16...\Q h4+ is the move.}

@ 17.Bf2
\mythoughts{17...\Q f6, targeting the long diagonal. 17...\Q e7 has no goal at all. With my move I am
also threatening ...\Q xf3.}

@ 18.Qe3
\mythoughts{White's queen on e3 is clumsy. His rook on a1 and knight on c3 are under attack from my
battery along the long diagonal. I want to activate my knight on h7: ...\N f8, ...\K h7 and ...\B h6
look attractive. 18...\N f8 is my move.}

@ 18...h4
\myreflection{Again, form without content. It may be true that I should activate my knight, but after
my line it is not activated at all. The idea of hitting with ...\B h6 also fails, because it hangs the
bishop. It is a pity that I make such simple mistakes as hanging a piece in my thinking process.
I feel I need to train my visualization.}

@ 19.h3
\mythoughts{I do not know how to proceed on the kingside. After 19...\B d7 20.O-O c4 21.bxc4 \N xc4
22.\Q f4, I have nothing better than trading the queens. 19...g5 is a good prophylactic move.}

@ 19...g5
\myreflection{The idea of \term{prophylaxis}\index{Strategy!Prophylaxis} is half right. Black has a knight
on e5 and must always keep an eye on f4: ...g5 is prophylaxis against f4. My queenside plan, however, is
the same as Kasparov's, and it is not the best, because it allows White to counterattack.}

@ 21.Kh1 | board
\mythoughts{21...\R ac8 and then ...c4 is my move: I can double my rooks on the c-file while White
struggles to find a move. 21...\N f8, heading for g6 and f4, is certainly a good move too.
Both should be good.}

@ 21...Nf8
\myreflection{Both moves are indeed fine. I fail to understand why the following line is better for White:
21...\R ac8 22.\R a2 c4?! (22...\N f8 is better) 23.bxc4 \N xc4 24.\Q a7.

When there are two good options, it is better to ask which one is absolutely necessary. Opening the
c-file is not: it brings Black no clear improvement, while improving the knight clearly does.}

@ 23.Bh2
\mythoughts{The f4 square is not available. If I play ...\B h6 and ...\N f4, the g-file will open.
I do not see how I can break through along the g-file, and besides, I am not prepared for it.
I still want to play ...\R ac8 with my intended ...c4.}

@ 24.Rb1
\mythoughts{I want to proceed with my plan. 24...\R ac8 is the move.}

@ 25.Qa7
\mythoughts{I did not see the queen move. I do not like 25...\N d3 26.\B xd6, because White suddenly
has a central passed pawn. 25...\N f4 covers my d6-pawn.

25...\N f4 26.bxc4 \R a8 27.\Q xb7 \R eb8 28.\Q xb8+ \R xb8 29.\R xb8+: I do not like my position.

25...\R a8 26.\Q e3 \R ac8, repeating the position, looks fine for me.}

@ 27.Rxb7 | board
\mythoughts{27...\N xe2 28.\B xe2 \B xh3 29.gxh3 \Q xc3 30.\R xf7: I am losing.

27...\N xe2 28.\B xe2 \Q xc3 29.\R xd7 \Q f6 30.\B xc4 \R xc4 31.\B xd6: I am losing.

I want to play 27...\R a8.}

@ 27...Nd2
\myreflection{My move 27...\R a8 is fine. In my line 25...\N f4 26.bxc4 \R a8? 27.\Q xb7 \R eb8 28.\Q xb8+
\R xb8 29.\R xb8+, I missed ...\N d2, as Kasparov played.

27...\N xe2 28.\B xe2 \Q xc3 29.\R xd7 \N e3 (instead of 29...\Q f6?) 30.\R g1 \Q d2 31.\R xf7 \B d4,
and the position is equal. Here I failed to realize that I did not need to protect the f7-pawn. My rooks
cover the back rank, and my bishop the vital g7 square. White's queen and rook cannot do further damage,
since his other pieces are far away. Black can, for example, keep harassing the queen with the rook:
...\R a8.}

@ 29.e5 | exercise
\mythoughts{29...\R a8 30.exf6 \R xa7 31.gxh3 \R xb7 32.fxg7: I am losing.

29...\Q f5 30.\R xd7 \Q xd7: I am winning.

29...\Q f5 30.e6 \B xe6 31.dxe6 \Q xe6: I cannot assess the position.

29...dxe5, threatening ...\R a8 again: I do not see how White can defend.

29...\R c5 also looks good.}

@ 29...Rxe5
\myreflection{29...\Q f5 30.e6 \B xe6 31.dxe6 \Q xe6 32.gxh3: White is winning.

29...\Q f5 30.e6 \R a8! (the same resource) 31.exd7 \R ed8 32.\Q b6 \N c4 33.\R f1 \N xb6 34.axb6
\B xc3: Black is winning.

29...dxe5 30.\R xd7 \N xg1 31.\Q xg1: the position is equal.

29...\R c5 30.\R f1 \Q f5 31.\R xd7 \B xe5 32.\B c2 \Q xc2 33.\B xe5 \Q f5 34.\B xd6 \R xc3 35.\R e7
\R cc8 36.\N d4 \Q f6 37.\R fe1 \N f2+ 38.\K g1 \N h3+ 39.\K h1 \N f2+: Black has perpetual check.

To be fair, I do not think I can calculate all these lines myself. One tip may help me not to miss my
moves: remember which resource comes up again and again in the position. Here I had already tried
...\R a8 in the previous moves. I should remember that it must be tried before I jump to conclusions.}

@ 30.Rxd7
\mythoughts{30...\N e4 31.\N xe4 \R xe4 32.fxe4 \N f2+ 33.\Q xf2 \Q xf2: I am winning.

30...\N e4 31.gxh3 \Q xf3+ 32.\R g2 \Q f1+ 33.\Q g1: White is winning.

30...\R c5 31.gxh3 \Q xf3+ 32.\R g2 \Q f1+ 33.\R g1: I have perpetual check.}

@ 31.Qxg1 | board
\mythoughts{31...\N xf3 32.gxf3 \Q xf3+ 33.\Q g2 \R xc3 34.\Q xf3 \R xf3 35.\B xe5 \B xe5:
Black has three pawns for one piece. I am no worse.}

@ 31...Ree8
\myreflection{31...\N xf3 32.gxf3 \Q xf3+ 33.\Q g2 \R e3 (my line 33...\R xc3 is a blunder because of
34.\R d8+) 34.\Q xf3 \R xf3 35.\N g1 \R fxc3: the position is indeed equal.}

@ 33.Ba4 | board
\mythoughts{33...\R xc3 34.\N xc3 leads nowhere, because the king is well defended.}

@ 33...Qd3
\myreflection{33...\Q d3 is a blunder. Here I failed to consider 33...\B xc3. Somehow this move was off my
radar, maybe because I value my bishop too much:
33...\B xc3 34.\B xe8 \R xe8 35.\N xc3 h3 36.\N e4 hxg2+ 37.\Q xg2 \N xf3 38.\R f6 \Q xe4 39.\R xf3 \Q xd5.

By the way, Nikitin does not mention this move in his book, and neither does Kasparov.}

@ 35.Rc6
\mythoughts{35...h3: I must try to expose the white king.}

@ 35...h3 | board

@ 36.Bg3
\mythoughts{36...\N xf3 37.gxf3 \Q xf3+ 38.\K h2 \R xe2+ 39.\N xe2 \Q xe2+ 40.\K xh3: I do not see
a perpetual check. 36...hxg2+ looks promising.}

\myreflection{36.\B g3 misses the win. 36.gxh3 \N xf3 37.\Q g2 \N xh2 38.\K xh2 \B xc3 39.\N xc3: the position
is clarified. White has an extra piece and is therefore winning.}

@ 39.Qf2 | board
\mythoughts{39...\Q h7+ 40.\B h2 \B e5 41.\N g3: I am clearly worse.
39...\Q f5, intending ...\Q h3+ with perpetual checks.}

\myreflection{White is in deep time trouble. To win the game, one must give him tough choices:
39...\B xc3, and 40.\N xc3? loses: 40...\R e1+ 41.\K g2 \R g1+ 42.\Q xg1 \N xg1 43.\R c8+ \K h7 44.d7 \Q xd7.

White must take with the rook: 39...\B xc3 40.\R xc3 \Q d1+ 41.\N g1 \N xg1 42.\Q xg1 \Q d5+ 43.\Q g2 \R h6+ 44.\B h2 \Q d1+ 45.\Q g1
\Q d5+, and Black has perpetual check.}

@ end
\subsection*{My Reflection}

It is a pity that I make simple mistakes, such as hanging pieces, in my planning, where no calculation
is involved. I must try to refute my own plans by asking myself: is the logic sound, have I missed
something, what else can go wrong? Hopefully, I will make fewer mistakes.

It is a very wild game, and I missed some moves. Asking what the typical resource in the position is
may help me avoid such misses.

I must do more visualization training with positions where many pieces are hanging.
